\documentclass[12pt]{article}

% House style preamble
\input{.house-style/preamble.tex}

% Update PDF metadata
\hypersetup{
  pdftitle={A metaphysician's systemic-functional turn: Thomasson's Rethinking Metaphysics and the autonomy of syntax},
  pdfkeywords={systemic functional linguistics, conceptual engineering, category and function, grammatical metaphor, autonomy of syntax, metaphilosophy}
}

% Working title -- provisional, confirm with Brett.
\title{A metaphysician's systemic-functional turn\\
{\large Amie Thomasson's \emph{Rethinking Metaphysics} and the autonomy of syntax}}
\author{Brett Reynolds \orcidlink{0000-0003-0073-7195}%
\thanks{Contact: \href{mailto:brett.reynolds@humber.ca}{brett.reynolds@humber.ca}}\\
Humber Polytechnic \& University of Toronto}
\date{\today}

\begin{document}
\maketitle

\begin{abstract}
\posscite{thomasson2025rethink} \emph{Rethinking Metaphysics} gives linguists a reason to care about a debate that might otherwise look internal to grammatical theory. Thomasson's constructive metaphilosophy draws on systemic functional linguistics to explain why nominalizations, modal vocabulary, and other philosophically loaded forms don't have to function by tracking objects or worldly joints. This review essay argues that the borrowing is important but architecturally risky. Thomasson is right that linguistic forms have different functions, and right that metaphysicians often project one tracking model across all discourse. But her conclusions don't require the Hallidayan placement of function inside grammar. A distribution-first, autonomy-of-syntax account can preserve the functional payoff of nominalization while keeping syntactic category, syntactic function, and pragmatic role distinct.
\end{abstract}

\noindent\textbf{Keywords:} systemic functional linguistics; conceptual engineering; category and function; grammatical metaphor; autonomy of syntax; metaphilosophy

\section{Introduction}
What's surprising about \posscite{thomasson2025rethink} \emph{Rethinking Metaphysics} is that the constructive half of the book depends on a Hallidayan way of asking what language is for. Thomasson turns to systemic functional linguistics to explain why easy ontological inferences are acceptable, why metaphysical disputes so often misfire, and how conceptual engineering should begin.

That should matter to linguists. A debate internal to grammar~-- whether functions organize syntax or are better kept distinct from syntactic description~-- has become part of a live metaphilosophical project. This paper's claim is correspondingly double. Thomasson is right that linguistic function matters for metaphysics. Her best results, though, don't require the systemic-functional placement of function inside grammar.

\section{Thomasson's systemic-functional turn}
Part I of \emph{Rethinking Metaphysics} is critical. Thomasson argues that familiar metametaphysical projects look for the wrong kind of uniform standard: joint-carving, grounding, truthmaking, causal access, or explanatory contribution. The constructive half changes the question. Instead of asking what all metaphysical discourse has to do, it asks what different bits of language are for, how they entered the language, and what speakers can do with them.

Systemic functional linguistics enters here. Thomasson says that SFL gives her a way to identify the functions of discourse and to explain why philosophically suspicious forms can still be in good order \citep[133--160]{thomasson2025rethink}. Her central examples are familiar from easy ontology: if the barn is red, then there's a property, redness; if there are electrons, then there's a number of electrons. A critic hears these transitions as cheap routes to abstract objects. Thomasson hears them as legitimate linguistic moves whose usefulness becomes clearer once we ask what nominalization lets speakers do.

Halliday's background claim isn't just that language is useful. In Halliday and Matthiessen's statement of the theory, language simultaneously construes experience, enacts social relations, and organizes discourse. The three modes are the ideational, interpersonal, and textual metafunctions. They're called \term{metafunctions} because, in the theory, function is internal to linguistic architecture: \enquote{the entire architecture of language is arranged along functional lines} \citep[30--31]{hallidayMatthiessen2014introduction}. That's the architectural claim at issue in this paper.

Thomasson uses the architecture most directly through grammatical metaphor. In Hallidayan terms, a meaning has a more congruent realization and may also have less congruent realizations, including nominalized forms. A clause can be reworked as a noun phrase; a process can be construed as a thing; a quality can support talk of a property. Thomasson takes this to explain why abstract-object vocabulary often misleads philosophers. If a noun entered the language through grammatical metaphor, we shouldn't expect it to work like a perceptual noun for a concrete object \citep[169--179]{thomasson2025rethink}.

For linguists, the news is concrete. SFL isn't a decorative citation in the book. It helps Thomasson move from negative deflationary arguments to a positive method for metaphysics as conceptual engineering. The method is reverse-engineering: identify a form, identify what earlier or more basic form it transforms, ask what needs that base form served, and then ask what the transformed form adds \citep[203--207]{thomasson2025rethink}. That's a serious use of a linguistic theory in philosophy, and it deserves a linguistic answer.

\section{Functions and their placement}
Start by keeping three uses of \term{function} apart. In the first use, familiar from \emph{The Cambridge Grammar of the English Language} and its descendants, a syntactic function is a relational position in constructional structure: subject, predicator, determiner, head, complement \citep{huddleston2002}. Function differs from lexical category. A determinative is a category; determiner is a function. The distinction matters because the same category can occur in different functions, and different categories can occur in the same function.

In the second use, \term{functional} is Hallidayan. Language is described in terms of what it enables speakers to do. Halliday's metafunctions~-- ideational, interpersonal, and textual~-- organize the grammar itself, rather than entering as later pragmatic uses of an independently described syntax. Thomasson adopts this picture when she turns to Halliday and Eggins for a systematic account of linguistic functions \citep[139--159]{thomasson2025rethink}.

In the third use, a function is a role in a pragmatic or teleological explanation. Thomasson has this sense in view when she asks what a stretch of discourse does for speakers, why it persists, and whether it should be retained, revised, or abandoned. It's also the sense that matters for conceptual engineering. A reverse-engineering project asks what a linguistic form is good for before deciding whether the form is philosophically suspect \citep[203--207]{thomasson2025rethink}.

Function matters. The question is where to put it. A distribution-first grammar can agree that nominalizations let speakers quantify, qualify, package information, and build theories. It can also agree that modal language serves social and normative ends. Those agreements leave open whether the ends are constitutive of syntactic category or syntactic function. The functions are real; the placement is optional.

Thomasson's borrowing becomes linguistically interesting at this point. For philosophers, systemic functional linguistics offers an attractive corrective to representational monism: many expressions earn their keep without naming or tracking an object. For an autonomist syntactician, the corrective comes bundled with an architectural choice. Halliday is right to insist that language serves multiple functions. On this view, his mistake is building those functions into the syntax.

\section{Grammatical metaphor and the category/function confound}
Grammatical metaphor is the pressure point. Thomasson is careful to block a bad philosophical reading of the term. A grammatical metaphor is a literal grammatical shift, as when \mention{red} in adjectival use supports the noun phrase \mention{the property of redness}, or when an event, process, or modal expression is recast in nominal form \citep[154]{thomasson2025rethink}. The resulting claims aren't make-believe. Mature language lets speakers recast material across grammatical forms.

Nominalization really does add capacities \citep[155--156]{thomasson2025rethink}. It lets speakers count, qualify, increase lexical density, package earlier claims, and construct layered theories. This lets Thomasson say that systemic functional linguistics gives her a new answer to a question left open by easy ontology: why would a language contain apparently redundant transitions from \mention{the barn is red} to \mention{the barn has the property of redness}? Such transitions are useful. They let speakers do work that unnominalized clauses don't do as well.

But usefulness isn't validity. Thomasson had already defended the easy inferences by appeal to introduction rules, redundancy, and ordinary linguistic competence before the systemic-functional apparatus enters the 2025 book \citep[34--39]{thomasson2025rethink}. She also says that, when she wrote \emph{Norms and Necessity}, she hadn't yet discovered systemic functional linguistics \citep[xiv]{thomasson2025rethink}; the modal-normativist program was already in place \citep{thomasson2020norms}. The SFL material helps explain why the transformations are worth having. By itself, it supplies no rule that licenses the inference.

This distinction matters because the SFL vocabulary makes a category shift look like a functional explanation. In the property case, the grammatical change is from adjective to noun. In CGEL terms, that's a change of category and construction. The functional payoff is a further fact about what the new construction lets speakers do. If we collapse the two, we risk treating a form's usefulness as if it explained its syntactic status, or treating a syntactic reclassification as if it by itself justified an ontological inference.

Thomasson's own later discussion pushes in the same direction. In Chapter 9, she says that the formal-functional analysis is incomplete on its own. Knowing that modal expressions or nominalizations serve broad formal functions doesn't yet distinguish the functions of different modal flavors or different grammatical metaphors. For that, she adds a needs-based, material analysis: identify the base form, then ask what needs or desires it serves and what later grammatical transformations add \citep[204--207]{thomasson2025rethink}. That's a retreat from pure grammatical architecture to a broader pragmatic genealogy.

Keep the functional turn. Thomasson's strongest conclusions survive a less Hallidayan grammar. She needs the claim that language has many functions, that nominalization expands what speakers can do, and that treating all nouns on the model of observational object tracking produces pseudo-problems. The stronger claim that those functions belong inside syntactic description is dispensable. Halliday gives her a powerful diagnostic vocabulary, but the inferential work is done by easy ontology plus framework-neutral facts about what nominalizations afford.

\section{What linguists should make of it}
First, linguistic theories travel. When they travel into philosophy, education, law, or cognitive science, their internal disputes often disappear from view. Thomasson's reader may come away thinking that Hallidayan grammar is simply what linguistics says once it takes function seriously. That would be a mistake. It's one powerful tradition within linguistics, not the neutral name for functional explanation.

Second, autonomist grammarians shouldn't answer Thomasson by denying the functional facts. Nominalization does let speakers repackage clauses as objects of quantification, modification, classification, and theoretical development. Modal vocabulary does do normative and social work. The better answer is to separate those facts from the claim that the grammar itself has to be organized by the functions. The autonomy of syntax isn't the claim that language has no uses. It's the claim that syntactic description shouldn't be made hostage to a prior functional taxonomy.

Third, Thomasson's examples expose a bad projection: a noun phrase needn't support the same inferences merely because it's a noun phrase. That point is compatible with a projectibility-first view of categories. The live question is what a category lets us infer, under which support conditions, with which failures and demotions, and across which domains. Some stable clusters are sustained by feedback or control processes; some aren't. Nothing in this argument requires treating homeostasis as the master explanation of every stable category.

For linguists, then, \emph{Rethinking Metaphysics} is both an opportunity and a warning. The opportunity is to show philosophers that detailed grammatical distinctions matter for live metaphysical questions. The warning is that, if linguists don't make the category/function distinction explicit, philosophers will inherit whichever architecture arrives bundled with the vocabulary they need. Thomasson needed a way to say that abstract-object talk can be useful without being observational object tracking. SFL gave her that way. A clearer grammar can keep the insight and remove the unnecessary commitment.

\section{Conclusion}
Thomasson's systemic-functional turn is a real event in the afterlife of linguistic theory. A major metaphysician has used SFL to rebuild the constructive part of metaphysics around the functions of discourse. Linguists shouldn't treat that as someone else's local philosophical borrowing.

The right response is neither dismissal nor conversion. Thomasson is right that metaphysicians have often mistaken one function of language for the function of language. She's also right that grammatical transformations such as nominalization explain why some abstract-object talk is ordinary, useful, and philosophically innocent. The caution is narrower: those results don't require the Hallidayan decision to put function inside syntax. They require a plural account of what linguistic forms do, plus a grammar that keeps category, function, and use in view without collapsing them.

\section*{Acknowledgements}

OpenAI Codex and Anthropic Claude Code served as drafting and editing aids during the preparation of this paper. I am responsible for all theoretical claims, arguments, errors, and interpretive choices.

\clearpage
\printbibliography

\end{document}
